\documentclass[12pt,a4paper]{article}

% ─── Packages ────────────────────────────────────────────────────────────────
\usepackage[utf8]{inputenc}
\usepackage[T1]{fontenc}
\usepackage{lmodern}
\usepackage[margin=2.5cm]{geometry}
\usepackage{graphicx}
\usepackage{hyperref}
\usepackage{booktabs}
\usepackage{listings}
\usepackage{xcolor}
\usepackage{amsmath}
\usepackage{enumitem}
\usepackage{setspace}
\usepackage{titlesec}
\usepackage{fancyhdr}
\usepackage{natbib}

% ─── Page style ──────────────────────────────────────────────────────────────
\pagestyle{fancy}
\fancyhf{}
\rhead{mobot\_core Research Report}
\lhead{\leftmark}
\cfoot{\thepage}
\setlength{\headheight}{14pt}

% ─── Hyperlink setup ─────────────────────────────────────────────────────────
\hypersetup{
    colorlinks=true,
    linkcolor=blue,
    citecolor=blue,
    urlcolor=blue,
    pdftitle={mobot\_core: Smartphone-Powered Robotics Research Report},
    pdfauthor={mobot\_core Project},
}

% ─── Code listing style ──────────────────────────────────────────────────────
\lstset{
    basicstyle=\ttfamily\small,
    keywordstyle=\color{blue}\bfseries,
    stringstyle=\color{red},
    commentstyle=\color{gray},
    numbers=left,
    numberstyle=\tiny\color{gray},
    breaklines=true,
    frame=single,
    captionpos=b,
    tabsize=4,
}

% ─── Title ───────────────────────────────────────────────────────────────────
\title{
    \vspace{2cm}
    \textbf{\Large mobot\_core} \\[0.5em]
    \large Smartphone-Powered Robotics: \\
    Research and Architecture Documentation \\[1em]
    \normalsize Inspired by OpenBot: Turning Smartphones into Robots \\
    \small \citep{openbot2021}
}
\author{mobot\_core Project}
\date{2024}

% ─── Document ────────────────────────────────────────────────────────────────
\begin{document}

\maketitle
\thispagestyle{empty}

\vfill
\begin{abstract}
    This report documents the research foundation of \texttt{mobot\_core}, a project exploring
    the use of commodity smartphones as the central computing and sensing platform for
    low-cost robotic systems. Inspired by the OpenBot project \citep{openbot2021}, which
    demonstrated that a smartphone-equipped robot body costing approximately \$50 can perform
    meaningful autonomous tasks, \texttt{mobot\_core} aims to document, understand, and
    eventually implement this approach.

    The report covers the background motivation for low-cost robotics, a detailed summary
    of the OpenBot paper, the proposed system architecture, hardware and software design,
    use cases, known limitations, and future implementation directions.
\end{abstract}
\vfill

\newpage
\tableofcontents
\newpage

% ─────────────────────────────────────────────────────────────────────────────
\section{Introduction}
\label{sec:introduction}
% ─────────────────────────────────────────────────────────────────────────────

Robotics research has historically been constrained by the high cost of capable hardware.
A standard research robot combining a compute unit, sensors, actuators, and a power system
can cost thousands to tens of thousands of dollars, placing serious robotics research beyond
the reach of most individuals and many institutions.

The central question that motivates \texttt{mobot\_core} is: \emph{Can we use the smartphone
that billions of people already carry in their pockets as the brain of a functional robot?}

Modern smartphones include remarkable sensing and computing capabilities:
\begin{itemize}
    \item High-resolution cameras (wide-angle, telephoto, depth)
    \item Inertial Measurement Units (accelerometer + gyroscope + magnetometer)
    \item GPS for outdoor localisation
    \item Powerful System-on-Chip (SoC) with CPU, GPU, and Neural Processing Unit (NPU)
    \item Wi-Fi, Bluetooth, and cellular connectivity
    \item A large battery (3000--6000 mAh)
    \item A mature operating system (Android) with rich software libraries
\end{itemize}

These capabilities are available in mid-range devices costing under \$100 — far less than
equivalent dedicated embedded hardware. The OpenBot project demonstrated that attaching a
smartphone to a simple \$50 robot body produces a system capable of person following and
autonomous navigation \citep{openbot2021}.

\texttt{mobot\_core} is a research and documentation project building on this foundation,
aiming to produce a well-understood architecture and a roadmap for a future implementation.

% ─────────────────────────────────────────────────────────────────────────────
\section{Background}
\label{sec:background}
% ─────────────────────────────────────────────────────────────────────────────

\subsection{The Cost Barrier in Robotics}

The high cost of capable robots is a well-known obstacle to democratising robotics research.
Commercially available research platforms include TurtleBot (\$500--\$1500), Universal Robots
arms (\$30,000+), and Boston Dynamics Spot (\$75,000+). Even the relatively affordable
Raspberry Pi-based robots require dedicated cameras, compute, and sensors that add up quickly.

\subsection{Smartphones as Embedded Computers}

The computational performance of smartphones has grown exponentially over the past decade.
A flagship smartphone SoC from 2023 (e.g., Snapdragon 8 Gen 2, Apple A16) delivers more
compute than a desktop computer from a decade ago. Importantly, even mid-range devices include
neural processing units capable of running machine learning inference at tens of milliseconds
per frame \citep{openbot2021}.

\subsection{TensorFlow Lite and Mobile AI}

The TensorFlow Lite framework \citep{tflite} allows deploying neural network models on
mobile and embedded devices with minimal dependencies. It supports hardware acceleration via:
\begin{itemize}
    \item GPU Delegate (OpenGL ES / Metal)
    \item NNAPI Delegate (Android Neural Networks API)
    \item Hexagon DSP Delegate (Qualcomm devices)
\end{itemize}
This makes on-device AI inference practical for real-time robotics workloads.

% ─────────────────────────────────────────────────────────────────────────────
\section{OpenBot Paper Summary}
\label{sec:paper_summary}
% ─────────────────────────────────────────────────────────────────────────────

\subsection{Research Problem}

M\"{u}ller and Koltun \citep{openbot2021} identified that the cost and complexity of existing
robotic platforms prevents broad participation in autonomous robotics research. Their goal was
to design a complete robot system — hardware and software — that any motivated individual
could build and use for serious research.

\subsection{Key Idea}

The key contribution is the observation that a commodity smartphone already contains all the
sensing and computing components needed for a capable robot brain. The remaining requirement
is a simple, cheap mechanical body for locomotion. This body can be built for approximately
\$50 using 3D-printed parts, off-the-shelf motors, and an Arduino Nano microcontroller.

\subsection{Main Contributions}

\begin{enumerate}
    \item A 3D-printable robot body design (chassis, motor mounts, phone holder)
    \item An Android application for teleoperation, sensor logging, and autonomous control
    \item A person-following model (MobileNet-SSD based) trained on collected data
    \item An autonomous navigation model (CNN, imitation learning from human demonstrations)
    \item A data collection pipeline for gathering training data with the robot
    \item Open-source release of all hardware designs, software, and trained models
\end{enumerate}

\subsection{Results}

The paper reports the following experimental results:

\begin{table}[h!]
\centering
\begin{tabular}{lll}
\toprule
\textbf{Task} & \textbf{Metric} & \textbf{Result} \\
\midrule
Person Following (indoor) & Success rate & $\approx$87\% \\
Person Following (outdoor) & Success rate & $\approx$78\% \\
Autonomous Navigation & Path completion rate & $\approx$72\% \\
Inference speed (mid-range phone) & Frames per second & 15--30 fps \\
Total hardware cost & USD & $\approx$\$50 \\
\bottomrule
\end{tabular}
\caption{OpenBot experimental results \citep{openbot2021}}
\end{table}

These results demonstrate that the smartphone-robot system is competitive with much more
expensive platforms for well-defined navigation and person-following tasks.

% ─────────────────────────────────────────────────────────────────────────────
\section{Proposed Architecture}
\label{sec:architecture}
% ─────────────────────────────────────────────────────────────────────────────

The \texttt{mobot\_core} architecture follows the same two-layer design as OpenBot, separating
high-level intelligence (smartphone) from low-level actuation (microcontroller).

\subsection{Smartphone Layer}

The smartphone handles:
\begin{itemize}
    \item \textbf{Perception}: Capturing camera frames and processing them with TFLite models.
    \item \textbf{State estimation}: Reading built-in IMU and GPS sensors.
    \item \textbf{Control}: Computing left/right wheel speed commands from perception outputs.
    \item \textbf{Communication}: Sending commands and receiving sensor feedback via USB serial.
    \item \textbf{Logging}: Recording all sensor and image data for training and analysis.
\end{itemize}

\subsection{Microcontroller Layer}

The Arduino Nano handles:
\begin{itemize}
    \item Parsing incoming serial commands from the smartphone.
    \item Setting PWM duty cycle on the L298N motor driver.
    \item Reading wheel encoder interrupts for odometry.
    \item Reading the HC-SR04 sonar sensor.
    \item Broadcasting sensor data to the smartphone.
\end{itemize}

\subsection{Communication Interface}

The smartphone communicates with the Arduino over USB serial (USB OTG) at 115,200 baud using
a simple text protocol:

\begin{lstlisting}[language=C, caption=Serial communication protocol]
// Phone -> Arduino
{l:150,r:150}        // Set left/right motor speeds (0-255)

// Arduino -> Phone
{l:148,r:152,s:45}   // Left speed, right speed, sonar (cm)
\end{lstlisting}

% ─────────────────────────────────────────────────────────────────────────────
\section{Hardware Design}
\label{sec:hardware}
% ─────────────────────────────────────────────────────────────────────────────

\subsection{Bill of Materials}

\begin{table}[h!]
\centering
\begin{tabular}{lll}
\toprule
\textbf{Component} & \textbf{Description} & \textbf{Cost (USD)} \\
\midrule
3D-printed chassis & PLA, differential drive frame & \$2--5 \\
TT gear motors $\times$ 2 & DC, 3--6V, ~150 RPM & \$4--8 \\
L298N motor driver & Dual H-bridge PWM & \$2--3 \\
Arduino Nano & ATmega328P microcontroller & \$3--5 \\
HC-SR04 sonar & Ultrasonic distance sensor & \$1--2 \\
Wheel encoders $\times$ 2 & Optical slotted disc + IR sensor & \$1--2 \\
2S LiPo battery & 7.4V, 1000--2000 mAh & \$8--12 \\
USB power bank & 5V/2A for phone charging & \$8--12 \\
Wheels and casters & 65mm wheels + caster ball & \$3--5 \\
USB OTG cable & Type-C or Micro-USB to Type-A & \$1--2 \\
Miscellaneous & Wires, fasteners, connectors & \$2--3 \\
\midrule
\textbf{Total} & & \textbf{\$35--59} \\
\bottomrule
\end{tabular}
\caption{Bill of materials for mobot\_core robot body}
\end{table}

\subsection{Power Architecture}

Two independent power rails are used:
\begin{itemize}
    \item \textbf{7.4V rail}: 2S LiPo $\rightarrow$ L298N $\rightarrow$ Motors; L298N 5V output
          $\rightarrow$ Arduino Nano.
    \item \textbf{5V rail}: USB power bank $\rightarrow$ Smartphone (via USB charging cable).
\end{itemize}

% ─────────────────────────────────────────────────────────────────────────────
\section{Software Design}
\label{sec:software}
% ─────────────────────────────────────────────────────────────────────────────

\subsection{Android Application Layers}

The Android application is structured in four layers:

\begin{enumerate}
    \item \textbf{Communication Layer}: USB serial driver, Bluetooth serial (optional).
    \item \textbf{Sensor Collection Layer}: CameraX, SensorManager, serial reader.
    \item \textbf{Perception Layer}: TensorFlow Lite inference engine, model management.
    \item \textbf{Control Layer}: Drive mode state machine, PID control, command generation.
\end{enumerate}

\subsection{Drive Modes}

\begin{table}[h!]
\centering
\begin{tabular}{ll}
\toprule
\textbf{Mode} & \textbf{Description} \\
\midrule
Manual & On-screen joystick or Bluetooth gamepad control \\
Person Following & AI detects person and steers towards them \\
Autonomous Navigation & CNN maps camera frames to wheel speeds \\
Data Collection & Human drives while app records camera and sensors \\
\bottomrule
\end{tabular}
\caption{Android application drive modes}
\end{table}

\subsection{AI Models}

\begin{itemize}
    \item \textbf{Person Detection}: MobileNet-SSD V2 (COCO pretrained), 300$\times$300 input,
          ~50--100 ms inference on CPU.
    \item \textbf{Navigation CNN}: Custom architecture trained via imitation learning, maps
          camera image to left/right speed command.
\end{itemize}

% ─────────────────────────────────────────────────────────────────────────────
\section{Use Cases}
\label{sec:use_cases}
% ─────────────────────────────────────────────────────────────────────────────

\texttt{mobot\_core} targets the following primary use cases:

\begin{description}
    \item[Education] Affordable hands-on robotics and AI learning platform for students.
    \item[Robotics research] Low-cost platform for imitation learning, visual navigation, and
          sensor fusion experiments.
    \item[Person following] Track and follow a specific person indoors and outdoors.
    \item[Autonomous navigation] Learn and replay navigation paths from human demonstrations.
    \item[Dataset collection] Collect labelled camera and sensor data for training AI models.
    \item[Social/assistive robotics] Companion or guide robot using the phone screen and speaker.
\end{description}

% ─────────────────────────────────────────────────────────────────────────────
\section{Limitations}
\label{sec:limitations}
% ─────────────────────────────────────────────────────────────────────────────

The following limitations are inherent to the low-cost commodity-hardware approach:

\begin{itemize}
    \item \textbf{Smartphone compatibility}: Performance varies significantly across devices.
          USB OTG support is not universal.
    \item \textbf{Battery life}: Approximately 30--60 minutes of motor operation; the phone
          battery drains quickly with continuous AI inference.
    \item \textbf{Sensor accuracy}: Consumer-grade IMU suffers from drift; sonar has limited
          angular coverage; GPS is unavailable indoors.
    \item \textbf{Real-time performance}: Android is not a real-time OS. End-to-end latency
          is 100--200 ms, limiting maximum safe operating speed.
    \item \textbf{Mechanical stability}: TT motors are not precisely matched; chassis can flex;
          performance degrades on uneven terrain.
    \item \textbf{Safety}: No hardware emergency stop; communication dropout risks.
    \item \textbf{Hardware variability}: Off-the-shelf components from different suppliers
          may have different characteristics.
\end{itemize}

% ─────────────────────────────────────────────────────────────────────────────
\section{Future Work}
\label{sec:future_work}
% ─────────────────────────────────────────────────────────────────────────────

The following development phases are planned for \texttt{mobot\_core}:

\begin{enumerate}
    \item \textbf{Arduino firmware}: Command parser, PWM motor control, PID speed control,
          encoder odometry, sonar reading, watchdog timeout.
    \item \textbf{Android application (manual)}: USB serial communication, camera preview,
          on-screen joystick, sensor dashboard, data logging.
    \item \textbf{Computer vision module}: TFLite person detection, person-following
          control law, GPU/NNAPI delegate support.
    \item \textbf{Autonomous navigation}: Imitation learning pipeline, CNN training,
          model deployment and evaluation.
    \item \textbf{Dataset collection pipeline}: Structured recording, export to standard formats.
    \item \textbf{SLAM support}: Visual-inertial odometry (ARCore), topological map building.
    \item \textbf{Simulation environment}: Gazebo or PyBullet model of the robot body.
    \item \textbf{Hardware evaluation}: Systematic testing of all components.
\end{enumerate}

% ─────────────────────────────────────────────────────────────────────────────
\section{Conclusion}
\label{sec:conclusion}
% ─────────────────────────────────────────────────────────────────────────────

This report has documented the research foundation of \texttt{mobot\_core}, a project
exploring smartphone-powered robotics. The central insight — that commodity smartphones
are capable sensing and computing platforms that can serve as robot brains — has been
validated empirically by the OpenBot project \citep{openbot2021}.

The key advantages of this approach are:
\begin{itemize}
    \item \textbf{Low cost}: A complete robot system for under \$60.
    \item \textbf{Capable AI}: On-device TFLite inference with GPU/NPU acceleration.
    \item \textbf{Rich sensing}: Camera, IMU, GPS, and wireless communication included.
    \item \textbf{Accessible}: Standard Android development tools and open-source ecosystem.
\end{itemize}

The \texttt{mobot\_core} project will build on this foundation to deliver a complete,
well-documented, open-source smartphone robotics platform. This documentation repository
serves as the knowledge base for all subsequent implementation work.

% ─────────────────────────────────────────────────────────────────────────────
\bibliographystyle{plainnat}
\bibliography{references}

\end{document}
